\section{Comparing empirical and theoretical distributions}

A central purpose of this chapter is to make the relationship between data and
models visible. Empirical summaries are computed from a finite sample. Theoretical
quantities belong to a probability law. The two levels can be compared, but they
should not be confused.

\subsection*{Theoretical CDF and empirical CDF}

Let \(F_X\) be the CDF of a theoretical model:
\[
F_X(x)=P(X\leq x).
\]
Let \(F_n\) be the empirical CDF of observed data:
\[
F_n(x)=\frac1n\sum_{i=1}^n\mathbf{1}_{\{x_i\leq x\}}.
\]
The function \(F_X\) gives probabilities according to the model. The function
\(F_n\) gives observed proportions according to the sample.

If the sample is generated from the model, then \(F_n(x)\) may be close to
\(F_X(x)\) for many values of \(x\), especially when \(n\) is large. But \(F_n\) is
still a step function built from data, while \(F_X\) is the theoretical CDF.

\subsection*{Theoretical mean and sample mean}

The theoretical mean is
\[
\mu=\E[X].
\]
The sample mean is
\[
\bar{x}=\frac1n\sum_{i=1}^n x_i.
\]
If the data are generated from a model with mean \(\mu\), then \(\bar{x}\) is an
empirical approximation to \(\mu\). For a small sample, the difference may be
large. For larger samples, the difference often becomes smaller under suitable
assumptions.

At this stage, the correct language is:
\begin{quote}
The sample mean is a data-based summary that may approximate the theoretical
mean.
\end{quote}
It is not correct to say that the sample mean is the theoretical mean.

\subsection*{Theoretical variance and sample variance}

The same distinction applies to variance. The model variance is
\[
\Var(X).
\]
The empirical variance is computed from data:
\[
v_n=\frac1n\sum_{i=1}^n(x_i-\bar{x})^2.
\]
The corrected sample variance
\[
s^2=\frac1{n-1}\sum_{i=1}^n(x_i-\bar{x})^2
\]
will become important when the aim is to estimate a population variance. But both
\(v_n\) and \(s^2\) are computed from the observed sample. They are not equal by
definition to the theoretical variance.

\subsection*{Histogram and density}

A histogram is built from data and depends on bin choices. A density is a
function from a probability model. A density-scaled histogram can be compared
with a density curve, but the comparison must be interpreted carefully.

If the histogram follows the shape of the density, then the sample is visually
consistent with the model. If the histogram is very different from the density,
then this may suggest that the model is inappropriate, the sample is small, the
binning is misleading, or the data contain structure not captured by the model.

A visual comparison is a starting point, not a final conclusion.

\subsection*{A guiding table}

\begin{center}
\begin{tabular}{lll}
\toprule
Question & Model object & Data object \\
\midrule
How much lies below \(x\)? & \(F_X(x)\) & \(F_n(x)\) \\
Where is the center? & \(\E[X]\) & \(\bar{x}\), median \\
How spread out is it? & \(\Var(X),\sigma_X\) & empirical variance, sample standard deviation, IQR \\
What is the shape? & density or PMF & histogram, frequency table \\
Where are probability positions? & theoretical quantiles & empirical quantiles \\
\bottomrule
\end{tabular}
\end{center}

This table summarizes the conceptual bridge built in this chapter. The next
chapters will study this bridge more deeply by treating sample summaries
as random variables.

\subsection*{Preparing for sampling distributions}

Once a sample is viewed as the realization of random variables
\[
X_1,
\ldots,X_n,
\]
statistics such as the sample mean also become random variables before the data
are observed:
\[
\bar{X}=\frac1n\sum_{i=1}^n X_i.
\]
After observation, \(\bar{X}\) takes the value \(\bar{x}\).

This is the next conceptual step. Descriptive statistics summarize the observed
sample. Sampling distributions describe how those summaries vary from sample to
sample before observation.

\begin{summarybox}
Empirical summaries connect data to probability models. They are not the same as
theoretical quantities, but they are designed to approximate, visualize, or
summarize the corresponding model concepts when data are generated from a model.
The next step is to study the randomness of the summaries themselves.
\end{summarybox}
